\chapter{Introduction}
\label{ch:intro}

\section{Motivation}

The open RISC-V instruction set architecture (ISA)~\cite{riscv-isa} has moved, in less than a
decade, from an academic teaching vehicle to a credible foundation for
commercial silicon. Its modular, royalty-free nature allows designers to tailor
a processor precisely to an application: a minimal microcontroller core for a
deeply embedded sensor node, or a richly extended core with custom accelerators
for signal processing and machine learning at the edge. This flexibility is
precisely what makes RISC-V attractive in domains where general-purpose
processors are either too power-hungry or too inflexible --- automotive
electronics, avionics and space payloads, industrial control, and battery-powered
edge artificial-intelligence (AI) devices.

These same domains, however, are also the domains in which \emph{dependability}
is not optional. A processor steering a vehicle, controlling an actuator, or
running a perception network must continue to produce correct results even in
the presence of transient and permanent faults. Such faults are no longer rare
curiosities confined to space applications. As CMOS technology has scaled into
the deep-nanometre regime, the electrical charge that distinguishes a logic
\texttt{1} from a logic \texttt{0} has shrunk dramatically, so that a single
energetic particle, a voltage glitch, or an ageing-induced threshold shift can
now flip a bit or momentarily corrupt a combinational result even at terrestrial
altitudes. Radiation-induced single-event effects (SEEs), supply-voltage noise,
electromigration, and bias-temperature instability all conspire to make modern
integrated circuits statistically more likely to err over their service life.

Within a modern embedded core, the \emph{vector accelerator} is a particularly
critical and particularly exposed component. Single-Instruction-Multiple-Data
(SIMD) accelerators exist precisely because they pack a large amount of
arithmetic into a small amount of control: wide datapaths of adders,
multipliers, shifters, and comparators that operate in lockstep on many data
elements at once. Consequently they tend to dominate both the silicon area and
the dynamic activity of the core. A fault that strikes the accelerator datapath
can silently corrupt every element of a vector result, and because the
accelerator is so heavily exercised, the probability that a fault is actually
\emph{activated} into an error is high. Protecting the vector accelerator is
therefore one of the highest-leverage decisions a designer can make when
hardening a RISC-V core for a critical application.

\section{The Klessydra vector accelerator}

This thesis works with the vector accelerator of \emph{Klessydra}, a family of
RISC-V cores developed for research into interleaved-multithreading and
SIMD acceleration. The Klessydra accelerator --- referred to throughout this
document as the \emph{DSP} unit --- implements a configurable SIMD datapath that
can be parameterised in its data width, the number of parallel SIMD lanes, the
size of the scratchpad memory, and the number of hardware threads it serves.
Internally it is organised around a small set of functional units: vector
adders, vector multipliers, vector shifters, vector comparators, and an
accumulation unit, fed by a mapping stage that distributes operands and gathers
results.

In its original form the accelerator was described as a largely monolithic block
of VHDL. While functionally complete, this structure is awkward for reliability
engineering: there is no clean boundary at which a redundancy mechanism can be
attached to, say, ``just the multiplier'' without disturbing the surrounding
logic. A central premise of this work is that \emph{fault tolerance is a
structural property as much as a functional one} --- to harden a design
selectively and to measure the cost of doing so, the design must first be
decomposed into well-defined, independently instantiable units.

\section{Problem statement}

Given the Klessydra vector accelerator, this thesis addresses three closely
related questions:

\begin{enumerate}
  \item \textbf{Structure.} How can the accelerator be re-engineered into a
        modular, hierarchical form in which each functional unit is cleanly
        isolated, so that redundancy can be inserted at arbitrary granularity ---
        from a single functional unit up to the entire accelerator?
  \item \textbf{Protection.} What reusable hardware redundancy mechanisms can be
        built on top of this modular structure, and what is the \emph{quantified
        cost} --- in area, power, and timing --- of applying them at different
        granularities and across different technology nodes?
  \item \textbf{Observability.} How can the hardened accelerator additionally be
        made \emph{observable in the field}, so that its health can be monitored
        at run time, and can this instrumentation be inserted automatically
        rather than by error-prone manual editing of the HDL?
\end{enumerate}

\section{Contributions}

In answering these questions, the thesis makes the following contributions.

\begin{itemize}
  \item \textbf{A modular re-engineering of the Klessydra DSP accelerator.}
        The monolithic datapath is restructured into a clean hierarchy in which
        adders, multipliers, shifters, comparators, and the accumulator are
        separate VHDL entities with explicit interfaces, multi-stage pipelined
        sub-blocks, and a mapping layer that is itself replication-aware.

  \item \textbf{Reusable DMR and TMR redundancy primitives.} Two small,
        technology-independent VHDL building blocks --- a Dual-Modular-Redundancy
        comparator with an error flag, and a Triple-Modular-Redundancy majority
        voter --- are designed and then wrapped around each functional unit and
        around the full accelerator, yielding a complete family of hardened
        configurations.

  \item \textbf{An automated IEEE~1687 (IJTAG) sensor-insertion tool.} A
        pragma-driven tool scans the VHDL hierarchy, inserts parameterisable
        on-chip sensors (activity and thermal monitors), and generates the
        Tessent scripts that extract the instrument connectivity and build the
        IJTAG access network with the necessary Segment Insertion Bits (SIBs),
        all without manual rewiring of the design.

  \item \textbf{A cross-node design-space exploration.} Every configuration is
        synthesised across three technology nodes spanning a mature, a
        mainstream, and an advanced process, and the resulting area, power, and
        timing data are reduced to a consistent set of comparisons, tables, and
        design guidelines.
\end{itemize}

\section{Thesis outline}

The remainder of this thesis is organised as follows.
Chapter~\ref{ch:background} reviews the technical background: the RISC-V ISA and
vector acceleration, the physics and modelling of faults in nanometre CMOS, the
classical redundancy techniques, and the IEEE~1149.1 / IEEE~1687 test and
monitoring standards.
Chapter~\ref{ch:architecture} describes the Klessydra DSP accelerator and the
modular re-engineering that underpins the rest of the work.
Chapter~\ref{ch:framework} presents the DMR and TMR fault-tolerance framework
and the two ways in which it is applied: per-functional-unit hardening and
replication-based redundancy.
Chapter~\ref{ch:ijtag} presents the automated IJTAG monitoring framework and the
sensor-insertion tool.
Chapter~\ref{ch:experiment} details the experimental setup: the synthesis flow,
the technology nodes, the configuration matrix, and the metrics.
Chapter~\ref{ch:results} reports and discusses the results of the design-space
exploration.
Chapter~\ref{ch:conclusion} concludes and outlines directions for future work.
